\documentclass[10pt,aspectratio=169]{beamer}

% ============================================================
% Módulo 3: Construcción de Escenarios Macro-Financieros
% Curso: Stress Testing en Sistemas de Pensiones de Capitalización Individual
% MACROPOL
% ============================================================

\usetheme{metropolis}
\usecolortheme{seahorse}
\usefonttheme{professionalfonts}

\usepackage[spanish,es-noquoting]{babel}
\usepackage[utf8]{inputenc}
\usepackage[T1]{fontenc}
\usepackage{amsmath,amssymb,bm}
\usepackage{booktabs}
\usepackage{array}
\usepackage{graphicx}
\usepackage{tikz}
\usepackage{listings}
\usepackage{xcolor}
\usetikzlibrary{arrows.meta,positioning,calc}

% ------------------------------------------------------------
% Colores
% ------------------------------------------------------------
\definecolor{MacropolBlue}{RGB}{0,67,122}
\definecolor{MacropolTeal}{RGB}{0,134,150}
\definecolor{MacropolGray}{RGB}{80,80,80}
\definecolor{SoftBlue}{RGB}{235,244,250}
\definecolor{SoftRed}{RGB}{252,238,238}
\definecolor{SoftGreen}{RGB}{235,248,240}
\definecolor{CodeBg}{RGB}{248,248,248}

\setbeamercolor{title}{fg=MacropolBlue}
\setbeamercolor{frametitle}{fg=MacropolBlue}
\setbeamercolor{structure}{fg=MacropolBlue}
\setbeamercolor{block title}{bg=MacropolBlue,fg=white}
\setbeamercolor{block body}{bg=SoftBlue,fg=black}
\setbeamertemplate{navigation symbols}{}
\setbeamertemplate{footline}[frame number]

% ------------------------------------------------------------
% R listings
% ------------------------------------------------------------
\lstdefinelanguage{R}{
  keywords={function,if,else,for,while,repeat,in,next,break,TRUE,FALSE,NULL,NA,NaN,Inf},
  otherkeywords={<-,<<-,$,~},
  sensitive=true,
  morecomment=[l]\#,
  morestring=[b]",
  morestring=[b]'
}

\lstset{
  language=R,
  basicstyle=\ttfamily\scriptsize,
  keywordstyle=\color{MacropolBlue}\bfseries,
  commentstyle=\color{MacropolGray}\itshape,
  stringstyle=\color{MacropolTeal},
  backgroundcolor=\color{CodeBg},
  frame=single,
  rulecolor=\color{MacropolGray!40},
  breaklines=true,
  showstringspaces=false,
  columns=fullflexible
}

% ------------------------------------------------------------
% Macros
% ------------------------------------------------------------
\newcommand{\baseline}{\textcolor{MacropolTeal}{\textbf{baseline}}}
\newcommand{\adverso}{\textcolor{red!70!black}{\textbf{adverso}}}
\newcommand{\R}{\texttt{R}}
\newcommand{\dd}{\mathrm{d}}

% ------------------------------------------------------------
% Metadata
% ------------------------------------------------------------
\title[Stress Testing Pensiones -- Módulo 3]{Módulo 3: Construcción de Escenarios Macro-Financieros}
\subtitle{ARIMA, VAR, Monte Carlo y diseño de choques para sistemas de pensiones}
\author{MACROPOL: Economía Aplicada y Políticas Públicas}
\institute{Curso: Stress Testing en Sistemas de Pensiones de Capitalización Individual}
\date{}

\begin{document}

% ============================================================
\begin{frame}
\titlepage
\end{frame}

% ============================================================
\begin{frame}{¿Dónde estamos en el curso?}
\begin{block}{Módulo 3 dentro del flujo de stress testing previsional}
El módulo transforma una \textbf{narrativa macroeconómica} en trayectorias cuantitativas de variables que luego alimentan el stress test del portafolio, el módulo actuarial y el reporte regulatorio.
\end{block}

\begin{center}
\begin{tikzpicture}[node distance=0.5cm, every node/.style={font=\scriptsize}]
\node[draw, rounded corners, fill=SoftBlue, minimum width=2.4cm, minimum height=0.8cm] (m1) {M1 Fundamentos};
\node[draw, rounded corners, fill=SoftBlue, right=of m1, minimum width=2.6cm, minimum height=0.8cm] (m2) {M2 Riesgos};
\node[draw, rounded corners, fill=SoftGreen, right=of m2, minimum width=3.0cm, minimum height=0.8cm] (m3) {\textbf{M3 Escenarios}};
\node[draw, rounded corners, fill=SoftBlue, right=of m3, minimum width=2.7cm, minimum height=0.8cm] (m4) {M4 Portafolio};
\node[draw, rounded corners, fill=SoftBlue, below=0.7cm of m4, minimum width=2.7cm, minimum height=0.8cm] (m5) {M5 Actuarial};
\node[draw, rounded corners, fill=SoftBlue, left=of m5, minimum width=3.0cm, minimum height=0.8cm] (m6) {M6 Reporte};

\draw[-{Latex}, thick] (m1) -- (m2);
\draw[-{Latex}, thick] (m2) -- (m3);
\draw[-{Latex}, thick] (m3) -- (m4);
\draw[-{Latex}, thick] (m4) -- (m5);
\draw[-{Latex}, thick] (m5) -- (m6);
\end{tikzpicture}
\end{center}

\textbf{Producto del módulo:} matriz de escenarios macro-financieros por horizonte, variable y severidad.
\end{frame}

% ============================================================
\begin{frame}{Objetivos de aprendizaje}
Al finalizar el módulo, los participantes podrán:

\begin{enumerate}
  \item Construir un escenario \baseline{} y uno \adverso{} para variables macro-financieras relevantes.
  \item Estimar modelos ARIMA para proyecciones univariadas.
  \item Estimar modelos VAR para capturar interdependencias macro-financieras.
  \item Generar trayectorias Monte Carlo y bandas de incertidumbre.
  \item Calibrar choques históricos e hipotéticos con trazabilidad regulatoria.
  \item Implementar un flujo reproducible en \R{}.
\end{enumerate}

% \vspace{0.4em}
% \begin{block}{Principio conductor}
% No buscamos ``predecir el futuro''. Buscamos construir trayectorias \textbf{coherentes, severas y plausibles} para evaluar la resiliencia del sistema previsional.
% \end{block}
\end{frame}

% ============================================================
\begin{frame}{Ruta gradual del módulo}
\begin{center}
\begin{tikzpicture}[node distance=0.35cm, every node/.style={font=\scriptsize}]
\node[draw, rounded corners, fill=SoftBlue, minimum width=2.5cm, minimum height=0.8cm] (n1) {1. Datos};
\node[draw, rounded corners, fill=SoftBlue, right=of n1, minimum width=2.5cm, minimum height=0.8cm] (n2) {2. ARIMA};
\node[draw, rounded corners, fill=SoftBlue, right=of n2, minimum width=2.5cm, minimum height=0.8cm] (n3) {3. VAR};
\node[draw, rounded corners, fill=SoftBlue, right=of n3, minimum width=2.5cm, minimum height=0.8cm] (n4) {4. Monte Carlo};
\node[draw, rounded corners, fill=SoftBlue, below=0.55cm of n4, minimum width=2.5cm, minimum height=0.8cm] (n5) {5. Choques};
\node[draw, rounded corners, fill=SoftGreen, left=of n5, minimum width=2.5cm, minimum height=0.8cm] (n6) {6. Taller R};

\draw[-{Latex}, thick] (n1) -- (n2);
\draw[-{Latex}, thick] (n2) -- (n3);
\draw[-{Latex}, thick] (n3) -- (n4);
\draw[-{Latex}, thick] (n4) -- (n5);
\draw[-{Latex}, thick] (n5) -- (n6);
\end{tikzpicture}
\end{center}

\begin{columns}
\column{0.48\textwidth}
\begin{block}{De menor a mayor complejidad}
\begin{itemize}
  \item Una variable: ARIMA.
  \item Varias variables: VAR.
  \item Muchas trayectorias: Monte Carlo.
  \item Narrativa adversa: choques calibrados.
\end{itemize}
\end{block}
\column{0.48\textwidth}
\begin{block}{De modelo a decisión}
\begin{itemize}
  \item Escenarios reproducibles.
  \item Supuestos auditables.
  \item Impactos comparables.
  \item Reglas de acción supervisora.
\end{itemize}
\end{block}
\end{columns}
\end{frame}

% ============================================================
\section{1. Contexto: de narrativa macro a escenario previsional}

% ============================================================
\begin{frame}{Contexto: ¿qué es un escenario macro-financiero?}
\begin{block}{Definición operativa}
Un escenario es una trayectoria conjunta de variables macroeconómicas y financieras bajo una narrativa específica, con horizonte, frecuencia, severidad y supuestos explícitos.
\end{block}

\begin{table}
\small
\begin{tabular}{p{3.2cm}p{4.2cm}p{4.2cm}}
\toprule
Componente & \baseline{} & \adverso{} \\
\midrule
Narrativa & Normalización gradual & Desaceleración + tasas altas + depreciación \\
Modelo & Proyección central & Proyección condicionada a choques \\
Supuestos & Política y mercado sin disrupciones & Persistencia de inflación y prima de riesgo \\
Resultado & Trayectoria esperada & Pérdida de valor y deterioro de aportes \\
\bottomrule
\end{tabular}
\end{table}
\end{frame}

% ============================================================
\begin{frame}{Canales e indicadores relevantes para pensiones de capitalización individual}
\begin{columns}
\column{0.52\textwidth}
\begin{block}{Canal financiero}
\begin{itemize}
  \item Tasa de interés local.
  \item Inflación y tasa real.
  \item Tipo de cambio.
  \item Spread soberano.
  \item Rentabilidad de bonos y renta variable.
\end{itemize}
\end{block}

\column{0.42\textwidth}
\begin{block}{Canal previsional}
\begin{itemize}
  \item Salario real.
  \item Empleo formal.
  \item Densidad de cotización.
  \item Aportes.
  \item Saldo acumulado.
\end{itemize}
\end{block}
\end{columns}

\vspace{0.4em}
\begin{center}
\begin{tikzpicture}[node distance=0.6cm, every node/.style={font=\scriptsize}]
\node[draw, rounded corners, fill=SoftRed, minimum width=3cm, minimum height=0.8cm] (shock) {Choque macro};
\node[draw, rounded corners, fill=SoftBlue, right=of shock, minimum width=3.4cm, minimum height=0.8cm] (risk) {Factores de riesgo};
\node[draw, rounded corners, fill=SoftGreen, right=of risk, minimum width=3.2cm, minimum height=0.8cm] (pension) {Balance previsional};
\draw[-{Latex}, thick] (shock) -- (risk);
\draw[-{Latex}, thick] (risk) -- (pension);
\end{tikzpicture}
\end{center}
\end{frame}

%------------------------------------------------
\begin{frame}{¿Qué son los factores de riesgo?}

En un sistema de capitalización individual, los \textbf{factores de riesgo} son variables que afectan:

\begin{itemize}
    \item el valor del portafolio previsional;
    \item el saldo acumulado del afiliado;
    \item la pensión esperada;
    \item la tasa de reemplazo.
\end{itemize}


\begin{block}{Definición}
Un factor de riesgo es toda variable que, al cambiar, modifica los activos, pasivos o beneficios previsionales.
\end{block}


\[
Resultado_t = f(F_1, F_2, \ldots, F_n)
\]

\begin{itemize}
    \item $F_i$: factores de riesgo
    \item $Resultado_t$: saldo, pensión o tasa de reemplazo
\end{itemize}


\textbf{Idea clave:} no se estresa directamente la pensión; se estresan los factores que la determinan.

\end{frame}

%------------------------------------------------
\begin{frame}{Tipos de factores de riesgo previsional}
\textbf{Módulo 2: Identificación y Medición de Riesgos Previsionales}

\begin{columns}[T]

\begin{column}{0.32\textwidth}
\begin{block}{Macroeconómicos}
\begin{itemize}
    \item PIB
    \item inflación
    \item tasas de interés
    \item desempleo
\end{itemize}
\end{block}
\end{column}

\begin{column}{0.32\textwidth}
\begin{block}{Financieros}
\begin{itemize}
    \item retornos
    \item volatilidad
    \item correlaciones
    \item riesgo de crédito
\end{itemize}
\end{block}
\end{column}

\begin{column}{0.32\textwidth}
\begin{block}{Actuariales}
\begin{itemize}
    \item longevidad
    \item edad de retiro
    \item densidad de cotización
    \item trayectoria salarial
\end{itemize}
\end{block}
\end{column}

\end{columns}


\begin{exampleblock}{Ejemplo de portafolio}
\[
R_p = \sum_{i=1}^{N} w_i R_i
\]

Los factores de riesgo son:
\[
R_i,\quad \sigma_i,\quad \rho_{ij}
\]
\end{exampleblock}

\end{frame}

%------------------------------------------------
\begin{frame}{Baseline vs. escenario adverso}
\textbf{Del factor de riesgo al impacto previsional}

\vspace{0.2cm}

\begin{table}
\centering
\small
\begin{tabular}{lcc}
\hline
\textbf{Factor} & \textbf{Baseline} & \textbf{Escenario adverso} \\
\hline
Retornos financieros & normales & caída de mercado \\
Tasa de interés & estable & shock de tasas \\
Inflación & controlada & aumento persistente \\
Longevidad & tendencia esperada & mayor esperanza de vida \\
Densidad de cotización & estable & caída por desempleo \\
\hline
\end{tabular}
\end{table}

\vspace{0.3cm}

\begin{block}{Intuición económica}
\begin{itemize}
    \item Shock financiero $\rightarrow$ menor saldo acumulado.
    \item Shock de longevidad $\rightarrow$ menor pensión mensual.
    \item Shock laboral $\rightarrow$ menor densidad de cotización.
    \item Shock combinado $\rightarrow$ deterioro de la tasa de reemplazo.
\end{itemize}
\end{block}

\vspace{0.2cm}

\textbf{Lectura regulatoria:} identificar qué factores explican la vulnerabilidad permite diseñar respuestas de política pública.

\end{frame}
%------------------------------------------------
% ============================================================
\begin{frame}{Modelo conceptual del escenario}
\begin{center}
\begin{tikzpicture}[node distance=0.6cm, every node/.style={font=\scriptsize}]
\node[draw, rounded corners, fill=SoftBlue, minimum width=3.0cm, minimum height=0.9cm] (data) {Datos históricos};
\node[draw, rounded corners, fill=SoftBlue, right=of data, minimum width=3.0cm, minimum height=0.9cm] (model) {Modelo estadístico};
\node[draw, rounded corners, fill=SoftRed, right=of model, minimum width=3.0cm, minimum height=0.9cm] (shock) {Choque adverso};
\node[draw, rounded corners, fill=SoftGreen, below=0.8cm of model, minimum width=3.6cm, minimum height=0.9cm] (scenario) {Trayectoria del escenario};
\node[draw, rounded corners, fill=SoftBlue, right=of scenario, minimum width=3.1cm, minimum height=0.9cm] (use) {Inputs M4--M5};

\draw[-{Latex}, thick] (data) -- (model);
\draw[-{Latex}, thick] (model) -- (shock);
\draw[-{Latex}, thick] (model) -- (scenario);
\draw[-{Latex}, thick] (shock) -- (scenario);
\draw[-{Latex}, thick] (scenario) -- (use);
\end{tikzpicture}
\end{center}

\begin{block}{Separación clave}
\textbf{Narrativa económica}: por qué ocurre el shock. \\
\textbf{Modelo cuantitativo}: cómo se propaga a las variables. \\
\textbf{Resultado}: cuánto cambia el balance previsional.
\end{block}
\end{frame}

% ============================================================
\begin{frame}{Datos: estructura mínima para el taller}
\begin{table}
\small
\begin{tabular}{lll}
\toprule
Variable & Frecuencia sugerida & Transformación \\
\midrule
PIB real (IMAE) & Mensual/Trimestral & Crecimiento interanual o trimestral anualizado \\
Inflación & Mensual / trimestral & Variación anual o trimestral \\
Tasa de política & Mensual / trimestral & Nivel o cambio en puntos básicos \\
Tipo de cambio & Mensual / trimestral & Depreciación porcentual \\
Spread soberano & Mensual / trimestral & Nivel o cambio en puntos básicos \\
Rentabilidad fondo & Mensual / trimestral & Retorno logarítmico \\
Salario real & Mensual / trimestral & Crecimiento real \\
\bottomrule
\end{tabular}
\end{table}

\begin{block}{Regla práctica}
El modelo más sofisticado no compensa datos mal transformados. La consistencia de frecuencia, unidades y signos es parte del control regulatorio.
\end{block}
\end{frame}

% ============================================================
\begin{frame}[fragile]{Implementación: plantilla de datos en R}
\begin{lstlisting}
# Paquetes sugeridos
library(readr)
library(dplyr)
library(lubridate)
library(tsibble)

datos <- read_csv("datos_macro_pensiones.csv") %>%
  mutate(fecha = ymd(fecha)) %>%
  arrange(fecha)

# Variables transformadas
datos_modelo <- datos %>%
  mutate(
    inflacion = 100 * (log(ipc) - lag(log(ipc), 12)),
    deprec_fx = 100 * (log(tc) - lag(log(tc), 12)),
    ret_fondo = 100 * (log(valor_cuota) - lag(log(valor_cuota))),
    salario_real = 100 * (log(salario_nominal/ipc) -
                          lag(log(salario_nominal/ipc), 12))
  ) %>%
  tidyr::drop_na()
\end{lstlisting}
\end{frame}

%------------------------------------------------
\begin{frame}{De shocks macro a impacto previsional: rol de ARIMA y VAR}

\textbf{Problema central del stress testing}

\vspace{0.2cm}

Los escenarios macroeconómicos (ej: recesión, crisis financiera) no impactan directamente las pensiones.

\begin{block}{Desafío}
¿Cómo transformar un \textbf{shock macro} en cambios cuantificables en:
\begin{itemize}
    \item retornos del portafolio,
    \item tasas de interés,
    \item salarios,
    \item densidad de cotización?
\end{itemize}
\end{block}

\vspace{0.3cm}

\textbf{Solución: Modelos de series de tiempo}
\end{frame}

\begin{frame}{De shocks macro a impacto previsional: rol de ARIMA y VAR}


\begin{columns}[T]

\begin{column}{0.48\textwidth}
\begin{block}{ARIMA}
\begin{itemize}
    \item Modela dinámica individual
    \item Captura persistencia y ciclos
    \item Útil para:
    \begin{itemize}
        \item inflación
        \item tasas de interés
    \end{itemize}
\end{itemize}
\end{block}
\end{column}

\begin{column}{0.48\textwidth}
\begin{block}{VAR}
\begin{itemize}
    \item Modela interacción entre variables
    \item Captura transmisión de shocks
    \item Útil para:
    \begin{itemize}
        \item PIB, tasas, mercados
        \item salarios y empleo
    \end{itemize}
\end{itemize}
\end{block}
\end{column}

\end{columns}

\vspace{0.3cm}

\begin{exampleblock}{Cadena de transmisión}
\[
\text{Shock macro} \rightarrow \text{Variables macro (VAR)} \rightarrow \text{Factores financieros/actuariales} \rightarrow \text{Pensión}
\]
\end{exampleblock}

\end{frame}

\begin{frame}{De shocks macro a impacto previsional: rol de ARIMA y VAR}

\begin{block}{Intuición económica}
\begin{itemize}
    \item Recesión $\rightarrow$ ↓ PIB $\rightarrow$ ↓ empleo $\rightarrow$ ↓ cotizaciones
    \item Crisis financiera $\rightarrow$ ↓ retornos $\rightarrow$ ↓ saldo acumulado
    \item Inflación persistente $\rightarrow$ ↓ tasas reales $\rightarrow$ ↓ rendimiento real
\end{itemize}
\end{block}

\vspace{0.2cm}

\textbf{Mensaje clave:} ARIMA y VAR permiten traducir narrativas macro en \textbf{trayectorias cuantitativas coherentes} para el stress testing.

\end{frame}
%------------------------------------------------

% ============================================================
\section{2. Modelos ARIMA}

% ============================================================
\begin{frame}{Contexto: ¿cuándo usar ARIMA en escenarios previsionales?}
\begin{block}{Uso principal}
ARIMA es el punto de entrada para construir trayectorias \baseline{} de una variable individual cuando queremos capturar persistencia, reversión y shocks idiosincráticos.
\end{block}

\begin{columns}
\column{0.48\textwidth}
\begin{block}{Útil para}
\begin{itemize}
  \item Inflación.
  \item Tasa de interés.
  \item Rentabilidad agregada.
  \item Salario real.
\end{itemize}
\end{block}

\column{0.48\textwidth}
\begin{block}{Limitación}
\begin{itemize}
  \item No impone coherencia entre variables.
  \item No explica contagio macro-financiero.
  \item Es más débil para escenarios condicionados.
\end{itemize}
\end{block}
\end{columns}
\end{frame}

% ============================================================
\begin{frame}{Modelo / Herramienta: ARIMA}
Sea $y_t$ una variable macro-financiera. Un modelo ARIMA$(p,d,q)$ se define como:

\[
\phi(L)(1-L)^d y_t = c + \theta(L)\varepsilon_t,
\quad \varepsilon_t \sim (0,\sigma^2)
\]

\begin{itemize}
  \item $p$: rezagos autorregresivos.
  \item $d$: diferenciación necesaria para estacionariedad.
  \item $q$: rezagos del error.
  \item $\phi(L)$ y $\theta(L)$: polinomios en el operador rezago.
\end{itemize}

\begin{block}{Intuición económica}
El pasado reciente contiene información sobre la velocidad de reversión, persistencia inflacionaria o inercia de tasas. ARIMA formaliza esa memoria.
\end{block}
\end{frame}

% ============================================================
\begin{frame}{ARIMA: flujo de estimación gradual}
\begin{enumerate}
  \item Graficar la serie y detectar quiebres, tendencia o estacionalidad.
  \item Transformar: log, diferencias, tasas de crecimiento.
  \item Evaluar estacionariedad.
  \item Seleccionar $(p,d,q)$ por AIC/BIC y criterio económico.
  \item Diagnosticar residuos.
  \item Proyectar \baseline{}.
  \item Condicionar o desplazar la trayectoria para un \adverso{} simple.
\end{enumerate}

\vspace{0.4em}
\begin{block}{Regla de enseñanza}
Primero se valida la historia económica de la serie; luego se valida la especificación estadística.
\end{block}
\end{frame}

% ============================================================
\begin{frame}[fragile]{Implementación: ARIMA para inflación}
\begin{lstlisting}
library(forecast)
library(ggplot2)

# Serie mensual de inflacion
infl_ts <- ts(datos_modelo$inflacion,
              start = c(2007, 1), frequency = 12)

# Seleccion automatica, pero revisada por el analista
fit_arima <- auto.arima(infl_ts,
                        seasonal = FALSE,
                        stepwise = FALSE,
                        approximation = FALSE)

summary(fit_arima)
checkresiduals(fit_arima)

# Proyeccion baseline
fc_base <- forecast(fit_arima, h = 8, level = c(50, 80, 95))
autoplot(fc_base)
\end{lstlisting}
\end{frame}

% ─────────────────────────────────────────────────────────────
\begin{frame}{De ARIMA a Escenario Adverso de Inflación}
  \footnotesize
\begin{block}{Regla general}
\[
\pi^{adv}_{t+h} = \hat{\pi}_{t+h} + k_h \cdot \sigma_{\varepsilon}
\]
\end{block}

\begin{itemize}
\item $\hat{\pi}_{t+h}$: pronóstico baseline (ARIMA)
\item $\sigma_{\varepsilon}$: desviación estándar de residuos
\item $k_h$: perfil dinámico del shock
\end{itemize}

\vspace{0.4cm}

\begin{block}{Perfil del shock ($k_h$)}
\begin{center}
\begin{tabular}{lcc}
\toprule
Horizonte & $k_h$ & Interpretación \\
\midrule
$t+1$ & 1.5 & shock máximo \\
$t+2$ & 1.5 & persistencia \\
$t+3$ & 1.5 & persistencia \\
$t+4$ & 1.0 & inicio reversión \\
$t+5$ & 0.8 & reversión \\
$t+6$ & 0.5 & disipación \\
$t+7+$ & 0.0 & normalización \\
\bottomrule
\end{tabular}
\end{center}
\end{block}
\end{frame}

%\vspace{0.3cm}
\begin{frame}{De ARIMA a Escenario Adverso de Inflación}
\begin{block}{Traducción a puntos porcentuales}
\[
Shock_{t+h} = k_h \cdot \sigma_{\varepsilon}
\]
\small
Ejemplo: si $\sigma = 0.8$, entonces shock máximo $= 1.5 \times 0.8 = 1.2$ p.p.
\end{block}

\vspace{0.3cm}

\begin{alertblock}{Supuesto clave}
La tasa nominal se mantiene constante para aislar el efecto inflacionario.
\end{alertblock}

\end{frame}

% ─────────────────────────────────────────────────────────────
\begin{frame}{Supuestos vs Resultados del Escenario}

\begin{columns}

\column{0.48\textwidth}
\begin{block}{Supuestos}
\begin{itemize}
\item Horizonte: 24 meses
\item Shock máximo: $1.5\sigma$
\item Persistencia: 3 meses
\item Reversión: gradual
\item Tasa nominal: baseline
\end{itemize}
\end{block}

\column{0.48\textwidth}
\begin{block}{Resultados}
\begin{itemize}
\item Inflación máxima
\item Shock en p.p.
\item Tasa real mínima
\item Meses con tasa real negativa
\end{itemize}
\end{block}

\end{columns}

\vspace{0.4cm}

\begin{alertblock}{Mensaje clave}
El modelo ARIMA genera el baseline. \\
El escenario adverso depende de supuestos explícitos del analista.
\end{alertblock}

\end{frame}



% ============================================================
\begin{frame}{Ejercicio ARIMA}
\begin{block}{Caso aplicado}
El regulador requiere un escenario de inflación para evaluar el retorno real de los fondos de pensiones durante los próximos 8 trimestres.
\end{block}

\begin{enumerate}
  \item Estimar un ARIMA para inflación.
  \item Generar la trayectoria \baseline{}.
  \item Construir un \adverso{} con inflación 1.5 desviaciones estándar por encima del baseline durante cuatro trimestres.
  \item Calcular tasa real aproximada:
  \[
  r^{real}_t \approx i_t - \pi_t
  \]
  \item Explicar el impacto esperado sobre bonos, salarios reales y aportes.
\end{enumerate}
\end{frame}

% ============================================================
\begin{frame}{Interpretación ARIMA para decisión regulatoria}
\begin{columns}
\column{0.48\textwidth}
\begin{block}{Qué mirar}
\begin{itemize}
  \item Persistencia del shock.
  \item Ancho de intervalos de predicción.
  \item Sesgo de residuos.
  \item Sensibilidad del retorno real.
\end{itemize}
\end{block}

\column{0.48\textwidth}
\begin{block}{Decisiones posibles}
\begin{itemize}
  \item Exigir escenarios de inflación más severos.
  \item Revisar supuestos de tasa real.
  \item Pedir análisis de sensibilidad de portafolio.
  \item Alinear horizonte con madurez de activos.
\end{itemize}
\end{block}
\end{columns}

\vspace{0.4em}
\begin{block}{Mensaje}
ARIMA produce una primera línea base. El stress test requiere luego consistencia multivariada.
\end{block}
\end{frame}

% ============================================================
\section{3. Modelos VAR}

% ============================================================
\begin{frame}{Contexto: por qué pasar de ARIMA a VAR}
\begin{block}{Motivación}
En pensiones, los shocks no llegan aislados: inflación, tasas, tipo de cambio, PIB y spreads se mueven conjuntamente. VAR captura esa propagación empírica.
\end{block}

\begin{center}
\begin{tikzpicture}[node distance=0.7cm, every node/.style={font=\scriptsize}]
\node[draw, rounded corners, fill=SoftRed, minimum width=2.2cm, minimum height=0.75cm] (fx) {FX};
\node[draw, rounded corners, fill=SoftRed, above right=of fx, minimum width=2.2cm, minimum height=0.75cm] (inf) {Inflación};
\node[draw, rounded corners, fill=SoftRed, below right=of fx, minimum width=2.2cm, minimum height=0.75cm] (spread) {Spread};
\node[draw, rounded corners, fill=SoftBlue, right=2.2cm of fx, minimum width=2.2cm, minimum height=0.75cm] (rate) {Tasa};
\node[draw, rounded corners, fill=SoftGreen, right=of rate, minimum width=2.2cm, minimum height=0.75cm] (port) {Portafolio};

\draw[-{Latex}, thick] (fx) -- (inf);
\draw[-{Latex}, thick] (fx) -- (spread);
\draw[-{Latex}, thick] (inf) -- (rate);
\draw[-{Latex}, thick] (spread) -- (rate);
\draw[-{Latex}, thick] (rate) -- (port);
\end{tikzpicture}
\end{center}
\end{frame}

% ============================================================
\begin{frame}{Modelo / Herramienta: VAR}
Sea $\bm y_t$ un vector de variables macro-financieras:

\[
\bm y_t =
\begin{bmatrix}
\Delta PIB_t \\
\pi_t \\
i_t \\
\Delta FX_t \\
spread_t
\end{bmatrix}
\]

Un VAR$(p)$ se escribe como:

\[
\bm y_t = \bm c + A_1 \bm y_{t-1} + \cdots + A_p \bm y_{t-p} + \bm u_t,
\quad \bm u_t \sim (0,\Sigma_u)
\]

\begin{block}{Intuición económica}
Cada variable responde a su propia historia y a la historia de las demás. El modelo estima cómo se transmite un shock de una variable al sistema.
\end{block}
\end{frame}

% ============================================================
\begin{frame}{VAR: decisiones técnicas críticas}
\begin{table}
\small
\begin{tabular}{p{3.2cm}p{7.6cm}}
\toprule
Decisión & Criterio aplicado \\
\midrule
Variables & Relevancia previsional y disponibilidad de datos \\
Transformación & Estacionariedad e interpretación económica \\
Número de rezagos & AIC/BIC, frecuencia y parsimonia \\
Ordenamiento & Supuesto de identificación de shocks \\
Estabilidad & Raíces dentro del círculo unitario \\
Diagnóstico & Autocorrelación, normalidad, heterocedasticidad \\
\bottomrule
\end{tabular}
\end{table}

\begin{block}{Regla regulatoria}
Todo VAR usado para stress testing debe documentar variables, transformaciones, rezagos, identificación y pruebas de estabilidad.
\end{block}
\end{frame}

% ============================================================
\begin{frame}[fragile]{Implementación: estimación VAR en R}
\begin{lstlisting}
library(vars)

Y <- datos_modelo %>%
  select(pib_crec, inflacion, tasa_policy,
         deprec_fx, spread_soberano) %>%
  as.matrix()

# Seleccion de rezagos
VARselect(Y, lag.max = 6, type = "const")

# Estimacion
fit_var <- VAR(Y, p = 2, type = "const")
summary(fit_var)

# Diagnosticos
serial.test(fit_var, lags.pt = 8, type = "PT.asymptotic")
arch.test(fit_var, lags.multi = 4)
roots(fit_var)  # estabilidad
\end{lstlisting}
\end{frame}

% ============================================================
\begin{frame}{Impulso-respuesta: leer la propagación del shock}
\[
IRF_{j,k}(h) = \frac{\partial y_{j,t+h}}{\partial u_{k,t}}
\]

\begin{block}{Ejemplo de lectura}
Un shock al tipo de cambio puede elevar inflación, inducir aumento de tasas, reducir crecimiento y afectar el valor de bonos. Esa secuencia es la narrativa cuantificada.
\end{block}

\begin{columns}
\column{0.50\textwidth}
\begin{itemize}
  \item $h=1$: impacto inmediato.
  \item $h=4$: persistencia anual.
  \item $h=8$: reversión o acumulación.
\end{itemize}

\column{0.45\textwidth}
\begin{itemize}
  \item Signo esperado.
  \item Magnitud económica.
  \item Duración del impacto.
  \item Coherencia con la narrativa.
\end{itemize}
\end{columns}
\end{frame}

% ============================================================
\begin{frame}[fragile]{Implementación: IRF y escenario condicionado}
\begin{lstlisting}
# Respuesta de variables ante shock cambiario
irf_fx <- irf(fit_var,
              impulse = "deprec_fx",
              response = c("inflacion", "tasa_policy",
                           "pib_crec", "spread_soberano"),
              n.ahead = 8,
              boot = TRUE,
              ci = 0.90)

plot(irf_fx)

# Pronostico baseline multivariado
fc_var <- predict(fit_var, n.ahead = 8, ci = 0.95)

# Extraer trayectoria central
base_var <- sapply(fc_var$fcst, function(x) x[, "fcst"])
\end{lstlisting}
\end{frame}

% ============================================================
\begin{frame}{Construcción de escenario adverso con VAR}
\begin{block}{Idea}
El escenario adverso puede generarse imponiendo un shock inicial y dejando que el VAR propague efectos dinámicos.
\end{block}

\[
\bm y_{t+1}^{adv} =
\hat{\bm c} + \hat A_1 \bm y_t + \cdots + \hat A_p \bm y_{t-p+1}
+ \bm s_{t+1}
\]

donde $\bm s_{t+1}$ es el vector de choques calibrados.

\vspace{0.6em}
\begin{table}
\small
\begin{tabular}{lcc}
\toprule
Variable & Shock inicial & Narrativa \\
\midrule
PIB & $-2\sigma$ & Desaceleración real \\
Inflación & $+1\sigma$ & Presión de precios \\
Tasa & $+150$ pb & Endurecimiento monetario \\
FX & $+2\sigma$ & Depreciación \\
Spread & $+200$ pb & Riesgo soberano \\
\bottomrule
\end{tabular}
\end{table}
\end{frame}

% ============================================================
\begin{frame}{Ejercicio VAR}
\begin{block}{Caso aplicado}
Diseñar un escenario adverso para el portafolio de pensiones ante depreciación cambiaria y aumento de spread soberano.
\end{block}

\begin{enumerate}
  \item Estimar un VAR con PIB, inflación, tasa interbancaria, tipo de cambio y spread.
  \item Construir una trayectoria adversa de 8 trimestres.
  \item Comparar contra el \baseline{}.
  \item Traducir los resultados a factores para el módulo 4: tasas, precios de bonos y rentabilidad esperada.
\end{enumerate}
\end{frame}

% ============================================================
\begin{frame}{Interpretación VAR para política pública}
\begin{columns}
\column{0.48\textwidth}
\begin{block}{Ventajas}
\begin{itemize}
  \item Coherencia entre variables.
  \item Propagación dinámica.
  \item Lectura de canales.
  \item Útil para sensibilidad supervisora.
\end{itemize}
\end{block}

\column{0.48\textwidth}
\begin{block}{Riesgos}
\begin{itemize}
  \item Relaciones inestables en crisis.
  \item Identificación sensible al ordenamiento.
  \item Muestras pequeñas.
  \item Extrapolación fuera de historia.
\end{itemize}
\end{block}
\end{columns}

\vspace{0.4em}
\begin{block}{Mensaje regulatorio}
El VAR no ``descubre'' el escenario adverso. Ayuda a hacerlo consistente, trazable y replicable.
\end{block}
\end{frame}

% ============================================================
\section{4. Simulación Monte Carlo}

% ============================================================
\begin{frame}{Contexto: de una trayectoria a muchas trayectorias}
\begin{block}{Problema}
Un solo escenario central oculta incertidumbre. Monte Carlo permite generar miles de trayectorias compatibles con el modelo y medir distribución de resultados.
\end{block}

\begin{center}
\begin{tikzpicture}[scale=0.9]
\draw[->] (0,0) -- (8,0) node[right]{horizonte};
\draw[->] (0,0) -- (0,3.2) node[above]{variable};
\foreach \i in {1,...,18}{
  \draw[MacropolBlue!35, thin] plot[smooth] coordinates {
    (0,1.4)
    (1,{1.4+rand*0.5})
    (2,{1.4+rand*0.7})
    (3,{1.5+rand*0.9})
    (4,{1.5+rand*1.1})
    (5,{1.6+rand*1.2})
    (6,{1.6+rand*1.3})
    (7,{1.6+rand*1.4})
  };
}
\draw[MacropolBlue, very thick] plot[smooth] coordinates {(0,1.4)(1,1.45)(2,1.5)(3,1.55)(4,1.6)(5,1.62)(6,1.65)(7,1.67)};
\end{tikzpicture}
\end{center}
\end{frame}

% ============================================================
\begin{frame}{Modelo / Herramienta: Monte Carlo sobre VAR}
Para un VAR estimado:

\[
\bm y_t = \hat{\bm c} + \hat A_1 \bm y_{t-1} + \cdots + \hat A_p \bm y_{t-p} + \bm u_t
\]

Simulamos:

\[
\bm u_t^{(m)} \sim \mathcal{N}(0,\hat \Sigma_u),
\quad m = 1,\ldots,M
\]

\[
\bm y_{t+h}^{(m)} =
\hat{\bm c} + \hat A_1 \bm y_{t+h-1}^{(m)} + \cdots + \hat A_p \bm y_{t+h-p}^{(m)}
+ \bm u_{t+h}^{(m)}
\]

\begin{block}{Intuición}
El modelo define la dinámica; Monte Carlo explora la incertidumbre de shocks futuros y su correlación.
\end{block}
\end{frame}

% ============================================================
\begin{frame}[fragile]{Implementación: simulación Monte Carlo simple}
\begin{lstlisting}
library(MASS)

B <- Bcoef(fit_var)          # coeficientes por ecuacion
Sigma <- summary(fit_var)$covres
K <- ncol(Y)
H <- 8
M <- 5000

# Funcion ilustrativa
sim_paths <- array(NA, dim = c(H, K, M))


for (m in 1:M) {
  y_lag <- tail(Y, 2)
  for (h in 1:H) {
    shock <- mvrnorm(1, mu = rep(0, K), Sigma = Sigma)
    y_new <- base_var[h, ] + shock
    sim_paths[h, , m] <- y_new
  }
}

# Percentiles por horizonte y variable
p05 <- apply(sim_paths, c(1, 2), quantile, probs = 0.05)
p50 <- apply(sim_paths, c(1, 2), quantile, probs = 0.50)
p95 <- apply(sim_paths, c(1, 2), quantile, probs = 0.95)
\end{lstlisting}
\end{frame}

% ============================================================
\begin{frame}{Cómo convertir simulaciones en escenarios}
\begin{table}
\small
\begin{tabular}{lll}
\toprule
Escenario & Regla cuantitativa & Uso \\
\midrule
Baseline & Percentil 50 & Planeación central \\
Adverso moderado & Percentil 10 o 5 & Sensibilidad supervisora \\
Adverso severo & Percentil 1 o combinación de colas & Resiliencia extrema \\
Reverse stress & Trayectoria que rompe umbral & Identificar vulnerabilidad crítica \\
\bottomrule
\end{tabular}
\end{table}

\begin{block}{Cuidado}
Los percentiles marginales no siempre forman una trayectoria económicamente coherente. Para stress testing, preferir trayectorias conjuntas simuladas o escenarios condicionados.
\end{block}
\end{frame}

% ============================================================
\begin{frame}{Ejercicio Monte Carlo}
\begin{block}{Caso aplicado}
El comité técnico quiere saber la probabilidad de que la tasa real sea negativa por cuatro trimestres y afecte la rentabilidad real del fondo.
\end{block}

\begin{enumerate}
  \item Simular 5,000 trayectorias de inflación y tasa nominal.
  \item Calcular $r^{real}_{t,h} \approx i_{t,h}-\pi_{t,h}$.
  \item Estimar:
  \[
  P(r^{real}_{t,h}<0 \text{ por 4 trimestres})
  \]
  \item Reportar percentiles de pérdida real acumulada.
  \item Traducir el resultado a una alerta de supervisión.
\end{enumerate}
\end{frame}

% ============================================================
\begin{frame}{Interpretación: de probabilidad a acción}
\begin{columns}
\column{0.48\textwidth}
\begin{block}{Indicadores técnicos}
\begin{itemize}
  \item Percentiles de variables.
  \item Probabilidad de umbral.
  \item Duración del shock.
  \item Correlación en estrés.
\end{itemize}
\end{block}

\column{0.48\textwidth}
\begin{block}{Acciones regulatorias}
\begin{itemize}
  \item Pedir análisis de liquidez.
  \item Revisar límites de concentración.
  \item Solicitar plan de comunicación.
  \item Aumentar frecuencia de monitoreo.
\end{itemize}
\end{block}
\end{columns}

\vspace{0.5em}
\begin{block}{Mensaje}
Monte Carlo no reemplaza al escenario narrativo. Lo complementa con distribución, probabilidad y sensibilidad.
\end{block}
\end{frame}

% ============================================================
\section{5. Choques históricos vs. hipotéticos}

% ============================================================
\begin{frame}{Contexto: dos familias de choques}
\begin{table}
\small
\begin{tabular}{p{3.0cm}p{4.2cm}p{4.2cm}}
\toprule
Tipo & Ventaja & Riesgo \\
\midrule
Histórico & Observable, comunicable, trazable & Puede no repetir estructura actual \\
Hipotético & Diseñado para vulnerabilidades actuales & Puede parecer arbitrario \\
Híbrido & Combina severidad observada y narrativa actual & Requiere mayor documentación \\
\bottomrule
\end{tabular}
\end{table}

\begin{block}{Principio}
El shock debe ser suficientemente severo para revelar vulnerabilidades, pero suficientemente plausible para guiar decisiones.
\end{block}
\end{frame}

% ============================================================
\begin{frame}{Choques históricos}
\begin{block}{Definición}
Reproducen movimientos observados en episodios pasados: crisis financiera global, pandemia, crisis cambiaria, alza abrupta de tasas o episodio de estrés soberano.
\end{block}

\begin{columns}
\column{0.48\textwidth}
\begin{block}{Cómo calibrar}
\begin{itemize}
  \item Identificar ventana de estrés.
  \item Medir cambios máximos acumulados.
  \item Escalar a frecuencia del modelo.
  \item Ajustar por condiciones actuales.
\end{itemize}
\end{block}

\column{0.48\textwidth}
\begin{block}{Aplicación previsional}
\begin{itemize}
  \item Caída de valor cuota.
  \item Aumento de tasa y pérdida en bonos.
  \item Menor salario real.
  \item Menor densidad de cotización.
\end{itemize}
\end{block}
\end{columns}
\end{frame}

% ============================================================
\begin{frame}[fragile]{Implementación: shock histórico}
\begin{lstlisting}
# Ventana historica de estres
shock_hist <- datos_modelo %>%
  filter(fecha >= ymd("2020-03-01"),
         fecha <= ymd("2020-12-31")) %>%
  summarise(
    shock_pib    = min(pib_crec, na.rm = TRUE),
    shock_infl   = max(inflacion, na.rm = TRUE),
    shock_tasa   = max(tasa_policy - first(tasa_policy), na.rm = TRUE),
    shock_fx     = max(deprec_fx, na.rm = TRUE),
    shock_spread = max(spread_soberano - first(spread_soberano),
                       na.rm = TRUE)
  )

shock_hist
\end{lstlisting}

\begin{block}{Lectura}
El episodio histórico se traduce a un vector de shocks. Luego se decide si se aplica directamente o se reescala según el balance de riesgos actual.
\end{block}
\end{frame}

% ============================================================
\begin{frame}{Choques hipotéticos}
\begin{block}{Definición}
Se diseñan desde una narrativa de riesgo prospectiva: política monetaria restrictiva prolongada, choque externo, salida de capitales o deterioro fiscal.
\end{block}

\begin{table}
\small
\begin{tabular}{lcc}
\toprule
Variable & Baseline & Adverso hipotético \\
\midrule
PIB real & 4.0\% & 0.5\% \\
Inflación & 4.5\% & 8.0\% \\
Tasa política & 6.0\% & 9.0\% \\
Depreciación & 3.0\% & 12.0\% \\
Spread soberano & 350 pb & 600 pb \\
Salario real & 2.0\% & -1.5\% \\
\bottomrule
\end{tabular}
\end{table}
\end{frame}

% ============================================================
\begin{frame}{Narrativa económica vs modelo cuantitativo}
\begin{columns}
\column{0.48\textwidth}
\begin{block}{Narrativa}
\begin{itemize}
  \item Choque externo eleva precios importados.
  \item Depreciación presiona inflación.
  \item Banco central sube tasa.
  \item Spread sube por aversión al riesgo.
  \item Crecimiento y empleo se debilitan.
\end{itemize}
\end{block}

\column{0.48\textwidth}
\begin{block}{Modelo}
\begin{itemize}
  \item VAR propaga relaciones dinámicas.
  \item ARIMA aporta baseline por variable.
  \item Monte Carlo cuantifica incertidumbre.
  \item Choques exógenos imponen severidad.
  \item Reglas de percentil documentan plausibilidad.
\end{itemize}
\end{block}
\end{columns}

\vspace{0.4em}
\begin{block}{Resultado}
Trayectoria trimestral coherente para alimentar modelos de portafolio y proyección de aportes.
\end{block}
\end{frame}

% ============================================================
\begin{frame}{Supuestos vs resultados}
\begin{table}
\small
\begin{tabular}{p{4.2cm}p{6.8cm}}
\toprule
Supuestos & Resultados \\
\midrule
Magnitud del shock inicial & Trayectoria proyectada por horizonte \\
Ordenamiento del VAR & Respuesta acumulada de variables \\
Ventana histórica elegida & Severidad observada y reescalada \\
Distribución de residuos & Percentiles Monte Carlo \\
Horizonte y frecuencia & Pérdidas temporales y recuperación \\
\bottomrule
\end{tabular}
\end{table}

\begin{block}{Disciplina de reporte}
Nunca presentar un resultado de stress testing sin separar explícitamente qué fue supuesto y qué fue generado por el modelo.
\end{block}
\end{frame}

% ============================================================
\begin{frame}{Ejercicio: comparar choque histórico e hipotético}
\begin{block}{Tarea}
Construir dos escenarios adversos de 8 trimestres para el mismo conjunto de variables.
\end{block}

\begin{enumerate}
  \item Escenario A: calibrado con la peor ventana histórica disponible.
  \item Escenario B: diseñado por narrativa hipotética de tasas altas y depreciación.
  \item Comparar severidad por variable.
  \item Identificar cuál es más dañino para:
  \begin{itemize}
    \item valor de bonos;
    \item retorno real;
    \item salario real;
    \item aportes esperados.
  \end{itemize}
  \item Elegir cuál debe ir al reporte regulatorio y justificar.
\end{enumerate}
\end{frame}

% ============================================================
\section{6. Taller práctico en R}

% ============================================================
\begin{frame}{Taller práctico: producto final}
\begin{block}{Entregable}
Un script reproducible en \R{} que genere una matriz de escenarios macro-financieros lista para ser usada en el stress testing del portafolio previsional.
\end{block}

\begin{table}
\small
\begin{tabular}{lllll}
\toprule
fecha & escenario & variable & valor & fuente \\
\midrule
2026T1 & baseline & inflación & 4.5 & ARIMA/VAR \\
2026T1 & adverso & inflación & 8.0 & VAR + shock \\
2026T1 & adverso & spread & 600 & shock hipotético \\
2026T1 & p05 & retorno real & -3.2 & Monte Carlo \\
\bottomrule
\end{tabular}
\end{table}
\end{frame}

% ============================================================
\begin{frame}{Estructura del script del taller}
\begin{enumerate}
  \item \texttt{00\_paquetes.R}: paquetes y funciones auxiliares.
  \item \texttt{01\_datos.R}: importación, limpieza y transformaciones.
  \item \texttt{02\_arima.R}: modelos univariados y baseline.
  \item \texttt{03\_var.R}: modelo multivariado e impulso-respuesta.
  \item \texttt{04\_montecarlo.R}: simulaciones y percentiles.
  \item \texttt{05\_choques.R}: escenarios histórico e hipotético.
  \item \texttt{06\_exportar.R}: tablas para módulos 4 y 5.
\end{enumerate}

\vspace{0.5em}
\begin{block}{Buenas prácticas}
Semilla fija, diccionario de variables, unidades explícitas, versiones de paquetes y exportación en formato largo.
\end{block}
\end{frame}

% ============================================================
\begin{frame}[fragile]{Implementación: objeto de escenarios}
\begin{lstlisting}
escenarios <- bind_rows(
  base_arima %>% mutate(escenario = "baseline_arima"),
  base_var   %>% mutate(escenario = "baseline_var"),
  adv_hist   %>% mutate(escenario = "adverso_historico"),
  adv_hip    %>% mutate(escenario = "adverso_hipotetico"),
  mc_p05     %>% mutate(escenario = "montecarlo_p05"),
  mc_p50     %>% mutate(escenario = "montecarlo_p50")
) %>%
  pivot_longer(cols = -c(fecha, escenario),
               names_to = "variable",
               values_to = "valor") %>%
  mutate(
    modulo_origen = "M3",
    fecha_generacion = Sys.Date()
  )

write_csv(escenarios, "outputs/escenarios_macrofinancieros.csv")
\end{lstlisting}
\end{frame}

% ============================================================
\begin{frame}{Guía de trabajo en grupos}
\begin{table}
\small
\begin{tabular}{p{2.7cm}p{4.2cm}p{4.2cm}}
\toprule
Grupo & Foco técnico & Pregunta de decisión \\
\midrule
A & ARIMA inflación y tasa & ¿El baseline es razonable? \\
B & VAR macro-financiero & ¿El shock se propaga con coherencia? \\
C & Monte Carlo & ¿Qué percentil debe usarse como adverso? \\
D & Choque histórico/hipotético & ¿Cuál escenario es más relevante para supervisión? \\
\bottomrule
\end{tabular}
\end{table}

\begin{block}{Discusión plenaria}
Cada grupo presenta una recomendación regulatoria, no solo un resultado estadístico.
\end{block}
\end{frame}

% ============================================================
\begin{frame}{Cómo enseñarlo en clase}
\begin{columns}
\column{0.48\textwidth}
\begin{block}{Secuencia docente}
\begin{enumerate}
  \item Intuición económica.
  \item Ecuación mínima.
  \item Código ejecutable.
  \item Diagnóstico.
  \item Decisión.
\end{enumerate}
\end{block}

\column{0.48\textwidth}
\begin{block}{Ritmo sugerido}
\begin{itemize}
  \item 20 min: concepto.
  \item 30 min: demostración en R.
  \item 40 min: práctica guiada.
  \item 20 min: interpretación regulatoria.
  \item 10 min: cierre y checklist.
\end{itemize}
\end{block}
\end{columns}

\vspace{0.4em}
\begin{block}{Principio pedagógico}
Cada técnica debe terminar en una tabla de decisión: variable, shock, severidad, canal previsional y acción supervisora.
\end{block}
\end{frame}

% ============================================================
\begin{frame}{Checklist regulatorio del escenario}
\begin{enumerate}
  \item ¿Está claramente separado el \baseline{} del \adverso{}?
  \item ¿La narrativa económica explica el origen del shock?
  \item ¿El modelo cuantitativo es replicable?
  \item ¿Los supuestos están documentados?
  \item ¿La severidad es extrema pero plausible?
  \item ¿Las variables son relevantes para pensiones?
  \item ¿El horizonte coincide con el uso del stress test?
  \item ¿Los resultados se exportan para portafolio y actuarial?
  \item ¿Existe análisis de sensibilidad?
  \item ¿La conclusión guía una decisión regulatoria?
\end{enumerate}
\end{frame}

% ============================================================
\begin{frame}{Cierre: mapa de decisión}
\begin{center}
\begin{tikzpicture}[node distance=0.65cm, every node/.style={font=\scriptsize}]
\node[draw, rounded corners, fill=SoftBlue, minimum width=3.0cm, minimum height=0.8cm] (arima) {ARIMA};
\node[draw, rounded corners, fill=SoftBlue, right=of arima, minimum width=3.0cm, minimum height=0.8cm] (var) {VAR};
\node[draw, rounded corners, fill=SoftBlue, right=of var, minimum width=3.0cm, minimum height=0.8cm] (mc) {Monte Carlo};
\node[draw, rounded corners, fill=SoftRed, below=0.85cm of var, minimum width=3.3cm, minimum height=0.8cm] (shock) {Choques};
\node[draw, rounded corners, fill=SoftGreen, below=0.85cm of shock, minimum width=4.0cm, minimum height=0.8cm] (esc) {Escenario macro-financiero};
\node[draw, rounded corners, fill=SoftBlue, below=0.85cm of esc, minimum width=4.0cm, minimum height=0.8cm] (decision) {Decisión regulatoria};

\draw[-{Latex}, thick] (arima) -- (var);
\draw[-{Latex}, thick] (var) -- (mc);
\draw[-{Latex}, thick] (var) -- (shock);
\draw[-{Latex}, thick] (mc) -- (shock);
\draw[-{Latex}, thick] (shock) -- (esc);
\draw[-{Latex}, thick] (esc) -- (decision);
\end{tikzpicture}
\end{center}

\begin{block}{Mensaje final}
El objetivo del módulo no es elegir el modelo ``más complejo'', sino construir escenarios que sean defendibles, reproducibles y útiles para proteger la suficiencia previsional.
\end{block}
\end{frame}

% ============================================================
\begin{frame}{Lectura aplicada sugerida}
\begin{itemize}
  \item Diseño de escenarios y calibración de choques: severidad, plausibilidad y objetividad.
  \item Uso de escenarios en stress testing: diferencias entre pruebas de estrés y métricas probabilísticas.
  \item Modelos macro-financieros para escenarios: ARMA/ARIMA, VAR, GARCH, BVAR y extensiones.
  \item Validación: robustez conceptual, sensibilidad y revisión independiente.
\end{itemize}

\vspace{0.6em}
\begin{block}{Para conectar con el módulo 4}
Exportar tasas, spreads, inflación, FX y retornos simulados como factores de riesgo del portafolio previsional.
\end{block}
\end{frame}

% ============================================================
\begin{frame}{Anexo: plantilla de discusión de escenarios}
\small
\begin{tabular}{p{2.7cm}p{8.3cm}}
\toprule
Elemento & Pregunta guía \\
\midrule
Narrativa & ¿Qué evento causa el escenario adverso? \\
Canal financiero & ¿Cómo afecta tasas, spreads, FX y retornos? \\
Canal previsional & ¿Cómo afecta aportes, salario real y saldo acumulado? \\
Modelo & ¿ARIMA, VAR, Monte Carlo o combinación? \\
Severidad & ¿Histórica, hipotética, percentil o híbrida? \\
Plausibilidad & ¿Por qué el escenario es defendible ante un comité? \\
Uso & ¿Qué decisión regulatoria se activa? \\
\bottomrule
\end{tabular}
\end{frame}

% ============================================================
\end{document}
